%-------------------------
% CV de Kevin Andres Gallardo Robles
% Version 3: diseno propio.
%   - Una columna, limpio y moderno, fondo blanco (legible e imprimible / ATS-friendly)
%   - Acento teal derivado del cian "Kinetic Obsidian" del portafolio
%   - Fira Sans (misma familia que el FiraMono de la v2, ya instalada)
% Compilar: lualatex cv_v3_custom.tex
%-------------------------

\documentclass[a4paper,10.5pt]{article}

\usepackage[
  top=1.3cm, bottom=1.3cm, left=1.6cm, right=1.6cm
]{geometry}

\usepackage[dvipsnames,table]{xcolor}
\usepackage{enumitem}
\usepackage{titlesec}
\usepackage{fontawesome5}
\usepackage[hidelinks]{hyperref}
\usepackage{tabularx}
\usepackage[spanish]{babel}

% Fira Sans como fuente principal (provista por el paquete fira de TeX Live).
% Bajo lualatex se carga vía fontspec con UTF-8 nativo, por eso NO se usa
% fontenc T1 (forzaba interpretación de 8 bits y rompía "·" y "—").
\usepackage[sfdefault]{FiraSans}

% --- Paleta -------------------------------------------------------------
% Acento inspirado en el cian #00F2FF del portafolio, oscurecido para
% mantener contraste sobre fondo blanco.
\definecolor{accent}{HTML}{0E7C86}
\definecolor{accentdark}{HTML}{0A5A62}
\definecolor{ink}{HTML}{1C2229}
\definecolor{muted}{HTML}{596573}
\definecolor{rule}{HTML}{D7DEE3}

\color{ink}
\hypersetup{urlcolor=accentdark, linkcolor=accentdark}

% --- Formato de secciones ----------------------------------------------
\titleformat{\section}
  {\large\bfseries\color{accent}}
  {}{0em}
  {\MakeUppercase}
  [{\color{rule}\titlerule[1pt]}]
\titlespacing{\section}{0pt}{10pt}{5pt}

% --- Listas -------------------------------------------------------------
\setlist[itemize]{
  leftmargin=1.1em, itemsep=1.2pt, topsep=2pt, parsep=0pt,
  label=\textcolor{accent}{\small\textbullet}
}

\setlength{\parindent}{0pt}
\pagestyle{empty}

% --- Helpers ------------------------------------------------------------
% Entrada con titulo a la izquierda y fecha a la derecha
\newcommand{\entry}[2]{%
  \noindent\begin{tabularx}{\textwidth}{@{}X@{}r@{}}
    \textbf{#1} & {\color{muted}\small #2}
  \end{tabularx}\par
}
% Subtitulo en itálica gris
\newcommand{\subentry}[1]{{\color{muted}\itshape #1}\par\vspace{2pt}}
% Pildora de stack
\newcommand{\stack}[1]{{\color{accentdark}\small\itshape #1}\par\vspace{1pt}}
% Item de la barra de habilidades
\newcommand{\skillrow}[2]{%
  \noindent\begin{tabularx}{\textwidth}{@{}p{3.1cm}@{}X@{}}
    \textbf{\color{ink}#1} & #2
  \end{tabularx}\par\vspace{2pt}
}

%-------------------------------------------------------------------------
\begin{document}

% ============ ENCABEZADO ============
{\fontsize{26}{30}\selectfont\bfseries\color{ink} Kevin Andrés Gallardo Robles}\\[3pt]
{\color{accent}\large Desarrollador Web Junior\, \textcolor{muted}{$|$}\, JavaScript \textcolor{muted}{$|$} React \textcolor{muted}{$|$} Node.js}\\[6pt]
{\color{rule}\rule{\textwidth}{1pt}}\\[5pt]
{\small
\faGlobe~\href{https://keevh.dev/}{keevh.dev}\quad
\faGithub~\href{https://github.com/keevh}{github.com/keevh}\quad
\faEnvelope~\href{mailto:andreskevin2606@gmail.com}{andreskevin2606@gmail.com}\quad
\faWhatsapp~3172172768
}

% ============ PERFIL ============
\section{Perfil Profesional}
\small
Desarrollador web junior y estudiante de \textbf{Ingeniería de Sistemas} en últimos semestres,
con experiencia práctica construyendo proyectos frontend y full stack con \textbf{JavaScript,
TypeScript, React, Next.js, Astro, Node.js, Express} y \textbf{Fastify}, además de bases de
datos relacionales y NoSQL. Oriento mi trabajo a interfaces web, consumo de APIs, lógica de
negocio, persistencia de datos, rendimiento y diseño responsive, con buen manejo de
Git/GitHub y disposición para aprender nuevas tecnologías según las necesidades del equipo.

% ============ HABILIDADES ============
\section{Habilidades Técnicas}
\small
\skillrow{Frontend}{JavaScript, TypeScript, React, Next.js, Astro, HTML, CSS, Tailwind CSS}
\skillrow{Backend}{Node.js, Express, Fastify, APIs REST, autenticación, validación de datos}
\skillrow{Bases de datos}{MongoDB, PostgreSQL, MySQL}
\skillrow{Herramientas}{Git, GitHub, Docker, npm, pnpm, despliegue web, documentación técnica}
\skillrow{Complementario}{Python, Java, IoT, MQTT, IA aplicada, Excel}

% ============ PROYECTOS ============
\section{Proyectos Destacados}

\entry{Lumen~\textcolor{muted}{\small— Demo de e-commerce con panel administrativo}}{\href{https://lumen.keevh.dev/}{\color{accentdark}\small lumen.keevh.dev}}
\stack{Next.js 16 · React 19 · TypeScript · Tailwind CSS · IndexedDB}
\begin{itemize}
  \small
  \item Tienda en español con catálogo, detalle de producto, carrito, checkout y panel administrativo para productos, descuentos, pedidos y configuración.
  \item Persistencia local con \textbf{IndexedDB} para simular un flujo completo de compra y administración sin backend externo.
  \item Arquitectura por capas (\textbf{Repository Pattern}) que desacopla la UI del almacenamiento y queda lista para evolucionar hacia una API real.
\end{itemize}
\vspace{4pt}

\entry{Trimly~\textcolor{muted}{\small— Acortador de URLs open source y autoalojable}}{\href{https://github.com/keevh/trimly}{\color{accentdark}\small github.com/keevh/trimly}}
\stack{TypeScript · Astro · React · Node.js · PostgreSQL · Docker}
\begin{itemize}
  \small
  \item App full-stack con \textbf{API REST} y \textbf{PostgreSQL}, desplegable con \textbf{Docker Compose}.
  \item Estadísticas públicas por enlace sin login, alias personalizados y expiración automática configurable.
  \item Borrado privado mediante token de gestión y \textbf{rate limiting por IP}; redirección separada del registro de analítica para no penalizar la latencia.
\end{itemize}
\vspace{4pt}

\entry{Otros proyectos}{\textcolor{muted}{\small juegos web desde cero}}
\begin{itemize}
  \small
  \item \textbf{\href{https://keevh.github.io/Alien-Monster-Hunter/}{\color{ink}Alien Monster Hunter}} \textcolor{accentdark}{\itshape(JavaScript, HTML5 Canvas, Vite, Web Audio API)}: shooter arcade con game loop, colisiones y oleadas construidos desde cero, audio sintetizado en tiempo real y modo de 1 y 2 jugadores.
  \item \textbf{\href{https://keevh.github.io/torres-de-hanoi/}{\color{ink}Torres de Hanoi}} \textcolor{accentdark}{\itshape(JavaScript, Web Audio API, localStorage, MockAPI)}: puzzle con drag and drop táctil, ranking por dificultad persistente y diseño responsive.
\end{itemize}

% ============ EXPERIENCIA ============
\section{Experiencia}
\entry{Desarrollo Web Independiente}{2024 -- Presente}
\subentry{Proyectos personales y a medida — Remoto}
\begin{itemize}
  \small
  \item Desarrollo aplicaciones web y prototipos funcionales que integran interfaz de usuario, lógica de negocio, consumo de APIs y persistencia de datos.
  \item Construyo proyectos de extremo a extremo: definición del problema, implementación, pruebas funcionales, despliegue y documentación técnica.
  \item Implemento funcionalidades de cliente y servidor: componentes reutilizables, rutas, validaciones, servicios, endpoints y manejo de datos.
  \item Aplico buenas prácticas de organización de código, separación de responsabilidades, diseño responsive y rendimiento web.
\end{itemize}

% ============ EDUCACION ============
\section{Educación}
\entry{Universidad Industrial de Santander (UIS)}{2019 -- En curso}
\subentry{Ingeniería de Sistemas — últimos semestres · Bucaramanga, Colombia}

\end{document}
